\section{Solutions for Chapter~\ref{ch:gapped}}

\paragraph{Exercise~\ref{ex:gapped-PF}.}
\phantomsection
\label{sol:ex:gapped-PF}
On the neutral basis $(f,s)$, with
$s=(o_0+o_1+o_2)/\sqrt3$, the rho matrix is
$R_0=\left(\begin{smallmatrix}1&\sqrt3\\\sqrt3&2\end{smallmatrix}\right)$.
Its characteristic polynomial is $x^2-3x-1$, giving $d$ and
$-d^{-1}$. On the color-orthogonal plane, $J$ vanishes and
rho is $-P$, where $P$ is the reflection of two colors. Hence
rho squares to one on that plane. Alpha is the order-three
rotation and $P\alpha P=\alpha^{-1}$.

On the neutral plane, $1+\alpha+\alpha^2=3$ and
$R_0^2=1+3R_0$. On the charged plane that alpha sum is zero
and $(-P)^2=1$. These two checks prove the full fusion relation
$R^2=1+(1+\alpha+\alpha^2)R$, as well as
$R\alpha=\alpha^{-1}R$. They identify the representation as
the two scalar fusion representations plus the two-dimensional
charged representation.



\paragraph{Exercise~\ref{ex:gapped-cubics}.}
\phantomsection
\label{sol:ex:gapped-cubics}
The ordered total probability is $p=3/(b^2+3)$ and the fixed
probability is $1-p=b^2/(b^2+3)$. Each ordered point has probability
$1/(b^2+3)$. The charges satisfy $\omega^{3j}=1$, so
\[
 \langle\Phi^3\rangle=p\,p^{-3/2}=p^{-1/2},\qquad
 \langle\Phi^\dagger\Phi E\rangle=p\,p^{-1}b/\sqrt3=b/\sqrt3.
\]
For the neutral cubic,
\[
 \langle E^3\rangle
 =\frac{b^3/\sqrt3-3\sqrt3/b}{b^2+3}
 =\frac{b^2-3}{\sqrt3 b}=-\frac1{\sqrt3},
\]
using $b^2+b=3$. Checking the multiplication identities at
the ordered and fixed points separately proves them as
identities of functions. For example, $1+gE$ equals
$1+b^2/3=p^{-1}$ on an ordered point and zero on the fixed
point, exactly the values of $\Phi^\dagger\Phi$.
Similarly $1-E/\sqrt3$ has values $b^2/3$ and $3/b^2$,
equal to $E^2$. Substituting $g=b/\sqrt3$ yields
$C_{EEE}=(g^2-1)/g$ and $p^{-1}=1+g^2$.
Finally $b^2=(7-\sqrt{13})/2$ gives
$p^{-1}=(13-\sqrt{13})/6$, recovering the charged cubic's
radical expression.



\paragraph{Exercises~\ref{ex:gapped-parent} and \ref{ex:gapped-short}.}
\phantomsection
\label{sol:ex:gapped-parent}
\label{sol:ex:gapped-short}
Invariance of a few global vectors constrains the operator only
on their span. It says nothing about its action on the local
range with arbitrary boundary data and arbitrary outside
environment. The stronger statement is that the open line
maps every interval module vector, with all allowed endpoint
and MPO boundary indices, into the complete local module
range. Closure under adjoints makes this range reducing, so
its orthogonal projector commutes with the closed line.

For the short Hamiltonian, one still needs equality of its
global zero space with the proposed MPS family and a positive
gap above that zero space uniformly in ring size. A different
longer parent can penalize directions left unpenalized by the
short terms. Its gap does not provide an operator comparison
bounding the short Hamiltonian from below by a positive
multiple of that parent. The four-site kernel count supplies
a concrete obstruction to the simplest such comparison.



\paragraph{Exercise~\ref{ex:gapped-filter}.}
\phantomsection
\label{sol:ex:gapped-filter}
With $B=Y_\rho-2$ and $B^\dagger B=5-Y_\rho$, commutation
with $H$ gives
\[
 E_B=\frac{5E-\langle HY_\rho\rangle}
                  {5-\langle Y_\rho\rangle}.
\]
Subtracting $E$ gives the covariance formula. The denominator
is at least $5-d>0$, and
$\operatorname{Var}(Y_\rho)\leq(d+d^{-1})^2/4=13/4$.
Cauchy--Schwarz therefore bounds the total energy difference
by
\[
 |E_B-E|\leq
 \frac{\sqrt{13}}{2(5-d)}
       \sqrt{\operatorname{Var}(H_L)}.
\]
For variance $O(L)$, division by $L$ gives $O(L^{-1/2})$.
A cat with nonzero weights on branches of densities $e_1\ne e_2$
has a contribution $w_1w_2(e_1-e_2)^2L^2$ to the variance.
For an explicit example, take $H_L=\sum_jZ_j$ on qubits and
$(|0\rangle^{\otimes L}+|1\rangle^{\otimes L})/\sqrt2$.
The branches have energies $L,-L$, so $\langle H_L\rangle=0$
and $\operatorname{Var}(H_L)=L^2$.


